\documentclass[12pt]{article}
\usepackage[T2A]{fontenc}
\usepackage[utf8]{inputenc}
\usepackage[russian]{babel}
\usepackage{amsmath}
\usepackage{amsfonts}
\usepackage{amssymb}
\usepackage{array}
\usepackage{geometry}
\geometry{margin=2.1cm}

\title{WonderFullyHE: математическая модель}
\date{}

\newcommand{\Enc}{\operatorname{Enc}}
\newcommand{\Dec}{\operatorname{Dec}}
\newcommand{\Encode}{\operatorname{Encode}}
\newcommand{\Decode}{\operatorname{Decode}}
\newcommand{\Resc}{\mathcal{R}}
\newcommand{\Relin}{\operatorname{Relin}}
\newcommand{\Rot}{\operatorname{Rot}}
\newcommand{\Switch}{\mathcal{S}}
\newcommand{\Lift}{\mathcal{L}}
\newcommand{\CtoS}{\Phi}
\newcommand{\StoC}{\Phi^{-1}}
\newcommand{\Red}{\operatorname{Red}}
\newcommand{\diag}{\operatorname{diag}}
\newcommand{\lvl}{\lambda}
\newcommand{\scl}{\delta}
\newcommand{\sz}{\nu}
\newcommand{\head}{\operatorname{head}}

\begin{document}
\maketitle

\section*{1. Входные объекты}

Библиотека работает с вектором
\[
    x=(x_0,\ldots,x_{m-1})\in\mathbb{C}^{m},
    \qquad
    1\leq m\leq S.
\]
Вещественные данные рассматриваются как частный случай:
\[
    x_i\in\mathbb{R}\subset\mathbb{C}.
\]

Параметры CKKS задаются набором
\[
    P=(N,\mathcal{Q},\Delta,S),
\]
где \(N\) --- степень полиномиального модуля, \(S=N/2\) --- число
доступных слотов, \(\Delta>0\) --- начальный масштаб, а
\[
    \mathcal{Q}=(q_0,\ldots,q_L)
\]
--- цепочка простых коэффициентных модулей.

На уровне \(\ell\) используется модуль
\[
    Q_\ell=\prod_{i=0}^{\ell}q_i,
    \qquad
    0\leq \ell\leq L.
\]
Рабочее кольцо:
\[
    R_{Q_\ell}=\mathbb{Z}_{Q_\ell}[X]/(X^N+1).
\]

Состояние шифртекста \(c\in R_{Q_\ell}^{2}\) описывается как
\[
    \operatorname{state}(c)=
    \bigl(\lvl(c),\scl(c),\sz(c),|\mathcal{Q}_{\ell}|\bigr),
\]
где \(\lvl(c)=\ell\) --- текущий уровень, \(\scl(c)\) --- масштаб,
\(\sz(c)\) --- число полиномов в шифртексте, а
\(|\mathcal{Q}_{\ell}|=\ell+1\) --- число модулей в текущей базе остатков.

Для режима обновления шифртекста используется профиль
\[
    N=16384,\qquad S=8192,\qquad \Delta=2^{40},
\]
\[
    (\log_2 q_0,\ldots,\log_2 q_L)
    =
    (60,40,40,40,40,40,40,40,60).
\]
В текущем проверочном контуре число логических слотов равно
\[
    m=16.
\]

\section*{2. Кодирование и шифрование}

Кодирование CKKS:
\[
    p=\Encode_{\Delta,\ell}(x)\in R_{Q_\ell}.
\]
Если уровень не указан, используется верхний уровень \(\ell=L\).

Шифрование:
\[
    c=\Enc_{pk}(p)\in R_{Q_\ell}^{2}.
\]

Расшифрование и декодирование:
\[
    \hat{x}=
    \Decode_{\scl(c)}
    \left(
        \Dec_{sk}(c)
    \right).
\]
Приближённость CKKS описывается равенством
\[
    \hat{x}=x+\varepsilon,
\]
где \(\varepsilon\) включает ошибку кодирования, шум шифрования, округления
при изменении масштаба и ошибки полиномиальных приближений.

\section*{3. Базовые операции над шифртекстами}

Пусть
\[
    c_x=\Enc(x),\qquad c_y=\Enc(y).
\]

Сложение:
\[
    c_{x+y}=c_x+c_y,
    \qquad
    \Dec(c_{x+y})\approx x+y.
\]
Вычитание:
\[
    c_{x-y}=c_x-c_y,
    \qquad
    \Dec(c_{x-y})\approx x-y.
\]

Если уровни различаются, шифртекст с большим уровнем переводится на уровень
второго слагаемого:
\[
    c_x'=\Switch_{\lvl(c_y)}(c_x).
\]
После этого проверяется относительное расхождение масштабов:
\[
    \rho=
    \frac{|\scl(c_x')-\scl(c_y)|}
    {\max(\scl(c_x'),\scl(c_y))}.
\]
В вычислительном слое используется ограничение
\[
    \rho\leq 10^{-3}.
\]

Умножение шифртекста на открытый вектор \(a\):
\[
    c_{a\odot x}=c_x\cdot\Encode(a),
    \qquad
    \Dec(c_{a\odot x})\approx a\odot x.
\]
Если после умножения выполняется изменение масштаба, то
\[
    c_{a\odot x}^{\,r}
    =
    \Resc(c_x\cdot\Encode(a)).
\]

Умножение двух шифртекстов:
\[
    c_{xy}
    =
    \Resc
    \left(
        \Relin(c_x\cdot c_y)
    \right),
\]
\[
    \Dec(c_{xy})\approx x\odot y.
\]
До изменения масштаба величина масштаба становится порядка
\[
    \scl(c_x)\scl(c_y).
\]
После изменения масштаба:
\[
    \lvl(\Resc(c))=\lvl(c)-1,
    \qquad
    \scl(\Resc(c))\approx\frac{\scl(c)}{q_{\lvl(c)}}.
\]
Здесь \(\mathcal{R}\) обозначает операцию изменения масштаба с удалением
одного модуля из текущей базы.

Ротация слотов:
\[
    \Rot_k(c_x)=\Enc(x^{(k)}),
    \qquad
    x^{(k)}_i=x_{(i+k)\bmod S}.
\]

\section*{4. Метрики ошибки}

Для точного результата \(y\) и расшифрованного результата \(\hat{y}\):
\[
    e_{\max}(y,\hat{y})
    =
    \max_{0\leq i<m}|y_i-\hat{y}_i|,
\]
\[
    e_{\operatorname{mean}}(y,\hat{y})
    =
    \frac{1}{m}\sum_{i=0}^{m-1}|y_i-\hat{y}_i|.
\]
Критерий корректности вычисления:
\[
    \operatorname{ok}(y,\hat{y},\tau)=
    \begin{cases}
        1,& e_{\max}(y,\hat{y})\leq\tau,\\
        0,& e_{\max}(y,\hat{y})>\tau.
    \end{cases}
\]
Для проверочного контура обновления шифртекста используется
\[
    \tau=2\cdot 10^{-5}.
\]

\section*{5. ABFT-инварианты}

Для полезной нагрузки
\[
    x=(x_0,\ldots,x_{m-1})
\]
задаётся контрольная сумма
\[
    \sigma(x)=\sum_{i=0}^{m-1}x_i.
\]
Расширенный вектор:
\[
    x^+=(x_0,\ldots,x_{m-1},\sigma(x)).
\]

Инварианты:
\[
    \sigma(x+y)=\sigma(x)+\sigma(y),
\]
\[
    \sigma(x-y)=\sigma(x)-\sigma(y),
\]
\[
    \sigma(x^{(k)})=\sigma(x),
\]
\[
    \sigma(x\odot y)=\sum_{i=0}^{m-1}x_iy_i.
\]
Проверка:
\[
    |\sigma_{\operatorname{obs}}-\sigma_{\operatorname{exp}}|
    \leq
    \tau_{\operatorname{abft}}.
\]
Порог \(\tau_{\operatorname{abft}}\) выбирается в соответствии с допустимой
численной ошибкой конкретного вычисления.

\section*{6. Линейные преобразования}

Пусть \(A\in\mathbb{C}^{S\times S}\). Для циклической диагонали \(k\):
\[
    d_k=(A_{0,k},A_{1,1+k},\ldots,A_{S-1,S-1+k}),
\]
где индексы берутся по модулю \(S\). Тогда
\[
    Ax=\sum_{k\in K}d_k\odot x^{(k)}.
\]
Над шифртекстом:
\[
    L_A(c_x)=
    \sum_{k\in K}\Encode(d_k)\cdot\Rot_k(c_x).
\]
После умножений на открытые диагонали члены приводятся к общему уровню и
масштабу, затем складываются.

Для плотных матриц используется разложение baby-step/giant-step:
\[
    b=\lceil\sqrt{|K|}\rceil,\qquad k=i+jb,\qquad 0\leq i<b.
\]
Предварительные ротации:
\[
    B_i=\Rot_i(c).
\]
Группы:
\[
    U_j=\sum_i\Encode(d_{i+jb})\cdot B_i.
\]
Итог:
\[
    L_A(c)=\sum_j\Rot_{jb}(U_j).
\]
Для \(m=16\) число необходимых ротационных ключей уменьшается с \(15\) до
\(6\).

\section*{7. Сумма слотов}

Для \(m=2^d\):
\[
    c^{(0)}=c,
\]
\[
    c^{(j+1)}
    =
    c^{(j)}+\Rot_{2^j}(c^{(j)}),
    \qquad
    j=0,\ldots,d-1.
\]
После \(d\) шагов:
\[
    \left(\Dec(c^{(d)})\right)_0
    \approx
    \sum_{i=0}^{m-1}x_i.
\]

\section*{8. Полиномиальные вычисления}

Для полинома
\[
    f(t)=\sum_{j=0}^{r}a_jt^j
\]
вычисление над шифртекстом задаётся как
\[
    f(c)=\sum_{j=0}^{r}a_jc^j.
\]
Каждое нетривиальное произведение шифртекстов расходует один уровень.
Перед сложением членов полинома выполняется перевод на общий уровень и
согласование масштабов.

\section*{9. Полиномиальная редукция}

В контуре обновления используется нечётный полином степени \(7\):
\[
    P_7(u)
    =
    u
    -6.579736267393u^3
    +12.98787880453u^5
    -12.20811674381u^7.
\]
Рабочий интервал:
\[
    |u|\leq B,
    \qquad
    B=2^{-10}.
\]
Опорная функция:
\[
    s(u)=\frac{\sin(2\pi u)}{2\pi}.
\]

Для шифртекста \(c_u\) вычисляются степени:
\[
    c_{u^2}=c_u^2,
    \qquad
    c_{u^3}=c_u\cdot c_{u^2},
\]
\[
    c_{u^5}=c_{u^3}\cdot c_{u^2},
    \qquad
    c_{u^7}=c_{u^5}\cdot c_{u^2}.
\]
После выравнивания уровней:
\[
    P_7(c_u)
    =
    c_u
    -6.579736267393c_{u^3}
    +12.98787880453c_{u^5}
    -12.20811674381c_{u^7}.
\]

\section*{10. Преобразования коэффициенты-слоты}

Для \(m\) логических слотов строится матрица канонического вложения
\[
    M\in\mathbb{C}^{m\times m}.
\]
Пусть
\[
    n=m,\qquad
    \zeta=\exp\left(\frac{2\pi i}{4n}\right),
    \qquad
    g_i=5^i\bmod 4n.
\]
Тогда
\[
    M_{ij}=\zeta^{g_i j},
    \qquad
    0\leq i,j<m.
\]

Переход из коэффициентного представления в слотовое:
\[
    z=Mx,
    \qquad
    c_z=\CtoS(c_x)=L_M(c_x).
\]
Обратный переход:
\[
    x=M^{-1}z,
    \qquad
    c_x=\StoC(c_z)=L_{M^{-1}}(c_z).
\]
Критерий обратимости на открытых данных:
\[
    M^{-1}Mx=x.
\]
Критерий на шифртексте:
\[
    \Dec\left(\StoC(\CtoS(c_x))\right)\approx x.
\]

\section*{11. Подъём уровня шифртекста}

Пусть шифртекст находится на уровне \(s<L\):
\[
    c\in R_{Q_s}^{2},
    \qquad
    Q_s=\prod_{i=0}^{s}q_i.
\]
Нужно представить те же полиномы в расширенной базе
\[
    Q_L=Q_s\prod_{i=s+1}^{L}q_i.
\]

Для каждого коэффициента \(a\in\mathbb{Z}_{Q_s}\) выбирается центрированный
представитель:
\[
    \tilde{a}=
    \begin{cases}
        a,&0\leq a\leq Q_s/2,\\
        a-Q_s,&Q_s/2<a<Q_s.
    \end{cases}
\]
Для новых модулей:
\[
    a_j=\tilde{a}\bmod q_j,
    \qquad
    j=s+1,\ldots,L.
\]
Расширенное представление:
\[
    a\mapsto
    (a\bmod q_0,\ldots,a\bmod q_s,a_{s+1},\ldots,a_L).
\]
Для полиномиальных коэффициентов это даёт отображение
\[
    \Lift:R_{Q_s}^{2}\rightarrow R_{Q_L}^{2}.
\]
Здесь \(\mathcal{L}\) обозначает подъём уровня шифртекста.
Масштаб сохраняется:
\[
    \scl(\Lift(c))=\scl(c),
\]
а уровень восстанавливается:
\[
    \lvl(\Lift(c))=L.
\]

\section*{12. Нормализация перед полиномиальной редукцией}

После перехода в слотовое представление:
\[
    z=Mx.
\]
Коэффициент нормализации:
\[
    \alpha=
    \begin{cases}
        1,&\max_i|z_i|\leq B/2,\\[4pt]
        \dfrac{B/2}{\max_i|z_i|},&\max_i|z_i|>B/2,
    \end{cases}
    \qquad
    B=2^{-10}.
\]
Тогда
\[
    \max_i|\alpha z_i|\leq 2^{-11}<B.
\]

Умножение декодируемого значения на \(\mu>0\) задаётся изменением масштаба:
\[
    \scl(c')=\frac{\scl(c)}{\mu}.
\]
Так как декодирование делит представление на текущий масштаб,
\[
    \Decode_{\scl(c')}\left(\Dec(c')\right)
    \approx
    \mu
    \Decode_{\scl(c)}\left(\Dec(c)\right).
\]
Уровень при этом не расходуется.

\section*{13. Обновление шифртекста}

Пусть входной шифртекст \(c_0\) находится на уровне \(\ell<L\). Обновление
задаётся последовательностью:
\[
    c_1=\Lift(c_0),
\]
\[
    c_2=\CtoS(c_1),
\]
\[
    c_3=\alpha c_2,
\]
\[
    c_4=P_7(c_3),
\]
\[
    c_5=\StoC(c_4),
\]
\[
    c_6=\alpha^{-1}c_5,
\]
\[
    c_7=\Lift(c_6).
\]
Результат:
\[
    c_{\operatorname{new}}=c_7.
\]

Критерий восстановления уровня:
\[
    \lvl(c_{\operatorname{new}})>\lvl(c_0).
\]
Критерий возможности продолжить вычисления:
\[
    \exists\,c_{\operatorname{next}}:
    c_{\operatorname{next}}
    =
    \Resc(c_{\operatorname{new}}\cdot\Encode(1)),
\]
\[
    \lvl(c_{\operatorname{next}})<\lvl(c_{\operatorname{new}}).
\]

В проверяемом режиме дополнительно считается
\[
    \hat{x}_{\operatorname{new}}
    =
    \head_m
    \left(
        \Decode
        \left(
            \Dec(c_{\operatorname{new}})
        \right)
    \right),
\]
и проверяется
\[
    e_{\max}(x,\hat{x}_{\operatorname{new}})
    \leq
    2\cdot 10^{-5}.
\]

\section*{14. Отчёт о вычислении}

Для каждого этапа \(t\) фиксируется
\[
    R_t=
    \left(
        n_t,
        \lvl_{\operatorname{in}},
        \lvl_{\operatorname{out}},
        \scl_{\operatorname{in}},
        \scl_{\operatorname{out}},
        e_{\max,t},
        T_t
    \right),
\]
где \(n_t\) --- название этапа, \(T_t\) --- время выполнения этапа.

Общий отчёт:
\[
    R=
    \left(
        R_1,\ldots,R_k,
        \alpha,
        \tau,
        \operatorname{preserve},
        \operatorname{restore}
    \right).
\]
Условия:
\[
    \operatorname{preserve}
    =
    \left[
        e_{\max}(x,\hat{x}_{\operatorname{new}})\leq\tau
    \right],
\]
\[
    \operatorname{restore}
    =
    \left[
        \lvl(c_{\operatorname{new}})>\lvl(c_0)
    \right].
\]

\section*{15. Общая схема}

\[
\begin{array}{c}
x\in\mathbb{C}^{m}
\\[4pt]
\downarrow\ \Encode_{\Delta}
\\[4pt]
p\in R_{Q_L}
\\[4pt]
\downarrow\ \Enc_{pk}
\\[4pt]
c\in R_{Q_L}^{2}
\\[4pt]
\downarrow\
\begin{array}{c}
\text{сложение, вычитание, умножение, ротации,}\\
\text{линейные преобразования, полиномиальные вычисления}
\end{array}
\\[10pt]
c_{\operatorname{eval}}
\\[4pt]
\downarrow\
\begin{array}{c}
e_{\max},e_{\operatorname{mean}}\\
\sigma(x)
\end{array}
\\[10pt]
c_{\operatorname{checked}}
\\[4pt]
\downarrow\
\begin{array}{c}
\Lift\\
\CtoS\\
\alpha\\
P_7\\
\StoC\\
\alpha^{-1}\\
\Lift
\end{array}
\\[10pt]
c_{\operatorname{new}}
\\[4pt]
\downarrow\ \Dec_{sk},\Decode
\\[4pt]
\hat{x}
\end{array}
\]

\section*{16. Формальный результат}

WonderFullyHE реализует отображение
\[
    \mathcal{F}:
    (x,P,\tau)
    \mapsto
    (c_{\operatorname{out}},\hat{x},R),
\]
где \(c_{\operatorname{out}}\) --- результирующий шифртекст,
\(\hat{x}\) --- приближённый расшифрованный результат, \(R\) --- отчёт
о вычислении.

Для обычного вычислительного пути:
\[
    e_{\max}(F(x),\hat{x})\leq\tau,
\]
где \(F\) --- функция, которую требуется вычислить над исходными данными.

Для обновления шифртекста:
\[
    \operatorname{preserve}=1
    \Longleftrightarrow
    e_{\max}(x,\hat{x}_{\operatorname{new}})\leq\tau,
\]
\[
    \operatorname{restore}=1
    \Longleftrightarrow
    \lvl(c_{\operatorname{new}})>\lvl(c_0).
\]

В текущем проверочном эксперименте:
\[
    \lvl(c_0)=6,
    \qquad
    \lvl(c_{\operatorname{new}})=7,
\]
и следующая операция с изменением масштаба переводит уровень
\[
    7\rightarrow 6.
\]

\end{document}
